\chapter*{Antes de começar}

Este livro é sobre uma máquina que pesa um quilo e meio, consome vinte watts, como uma lâmpada fraca, e mesmo assim reconhece rostos, aprende línguas e apanha uma bola no ar. É sobre o seu cérebro. Vamos olhar para ele como um engenheiro olha para uma placa desconhecida: que peças são essas, como estão ligadas, que sinais correm pelos fios e o que esse circuito calcula.

Você não vai precisar de fórmulas nem de conhecimentos de biologia. Vai precisar de curiosidade e de disposição para imaginar um fio, uma lâmpada, um interruptor. Todo o resto vamos montar pelo caminho.

O livro tem uma segunda versão, completa --- com equações, modelos, exercícios e referências aos experimentos. Os capítulos das duas versões seguem a mesma ordem e têm os mesmos números. Se algum capítulo prender a sua atenção, abra a versão completa dele: lá você vai encontrar de onde veio cada afirmação e como verificá-la por conta própria. No site do livro, cada capítulo pode ser alternado entre ``Resumo'' e ``Completo'' com um único botão.

Um aviso. Sobre o cérebro, está longe de se saber tudo, e vou dizer com honestidade onde termina o que foi medido e começa o palpite. Isso não é um defeito do livro, e sim o que ele tem de mais interessante: muitos capítulos terminam com perguntas cuja resposta, por enquanto, ninguém conhece.
